% chapters/07/checklist.tex
\bagian{Checklist Pencapaian Kompetensi}

Bab ini menuntaskan \SUBCPMK{1.4} dan menutup \CPMK{1}. Centang setiap butir
berikut hanya apabila Anda benar-benar telah melakukannya atau dapat
menjelaskannya.

\begin{checklistbox}
\begin{checklist}
  \item {\bfseries (Indikator a)} Saya dapat menyusun benda permainan sebagai
        kelas turunan \code{pygame.sprite.Sprite} dengan anggota
        \code{image} dan \code{rect} yang benar.
  \item Saya dapat mengelola banyak benda dengan \code{pygame.sprite.Group},
        termasuk memanggil \code{update} dan \code{draw} untuk seluruhnya.
  \item Saya dapat memilih di antara \code{colliderect},
        \code{spritecollide}, dan \code{groupcollide} sesuai keperluan.
  \item Saya dapat menjelaskan kapan persegi pembatas tidak memadai dan
        bagaimana \istilah{mask} dipakai sebagai gantinya.
  \item Saya dapat menjelaskan pertimbangan biaya perhitungan antara kedua
        metode itu.
  \item {\bfseries (Indikator b)} Saya dapat menulis penyelesaian tabrakan
        sumbu demi sumbu, mendahulukan sumbu mendatar.
  \item Program saya membuat pemain dapat mendarat di atas bongkahan sekaligus
        terhalang oleh sisinya.
  \item Saya dapat menjelaskan mengapa kecepatan pada sumbu yang diselesaikan
        perlu dinolkan.
  \item Penghitung tumpang tindih pada program saya bernilai nol pada seluruh
        laju bingkai yang saya uji.
  \item {\bfseries (Indikator c)} Saya dapat menyebutkan tiga gejala tabrakan
        yang salah beserta penyebabnya.
  \item Saya dapat menghitung mengapa benda menembus bongkahan yang tipis pada
        kecepatan tertentu.
  \item Saya dapat menjelaskan mengapa benda terangkat ke atas bongkahan
        ketika kedua sumbu diperiksa sekaligus.
  \item Saya telah merusak program dengan sengaja, mengamati gejalanya, lalu
        memperbaikinya kembali.
  \item Catatan praktikum saya memuat ketiga gejala itu beserta perbaikannya.
\end{checklist}
\end{checklistbox}
